\section{Related Work}\label{sec:related_work}

\subsection{Physics-based motion imitation}

DeepMimic demonstrated reference-guided reinforcement learning for simulated
character skills~\cite{deepmimic}. Perpetual Humanoid Control developed robust
humanoid imitation for real-time simulated avatars~\cite{phc}. These approaches
motivate learning a feedback controller rather than directly replaying a
kinematic trajectory. Subsequent single-body trackers scaled this
feedback-control approach in coverage and deployability: PULSE distills a
low-dimensional, high-coverage motion representation for general humanoid
control~\cite{pulse}; OmniH2O and ExBody2 target dexterous, deployable
whole-body teleoperation and expressive control on physical
humanoids~\cite{omnih2o,exbody2}; GMT uses a motion mixture-of-experts to
track a broad motion distribution with one policy~\cite{gmt}. Each of these
typically targets a fixed robot platform or morphology while scaling the
motion distribution; we instead fix the tracking task and scale the body,
holding 128 configurations under one
policy. We implement our experiments in
ProtoMotions~\cite{protomotions} with Isaac Gym~\cite{isaacgym}, using PPO
optimization~\cite{ppo}. Our contribution is the multi-morphology data and
control workflow and its evaluation, rather than a replacement for these
simulation or learning frameworks.

\subsection{Morphology-conditioned control and retargeting}

Body-shape variation has already been studied in physics-based control:
Won and Lee learned a parametric controller that accommodates variation in
character height, weight, and proportions~\cite{won2019}. A separate line of
work learns policies that transfer across changes in agent \emph{topology} --
limb count and skeletal structure -- through shared modular
sub-policies~\cite{smp2020}, a different axis from the continuous
anthropometric variation studied here. Morphology conditioning is therefore
an established idea; our contribution is not the concept but a system that
examines one controller over a large set of paired motions and bodies, with
shape identity maintained throughout data generation, simulation, and
evaluation.

Table~\ref{tab:related_work_comparison} summarizes how this positions against
the work discussed above. No shared benchmark or task spans these systems, so
the comparison is qualitative rather than a head-to-head evaluation; to the
best of our knowledge, none of them trains and evaluates one policy across a
paired, multi-shape motion corpus at anywhere near this scale.

\begin{table}[t]
\centering
\small
\caption{Qualitative positioning against the physics-based imitation and
morphology-conditioned control work discussed above. ``Body configs'' counts
distinct trained body shapes evaluated under one policy, not motion coverage;
entries reflect how each work is characterized in this section, not a full
independent audit of those papers.}\label{tab:related_work_comparison}
\begin{tabularx}{\linewidth}{@{}lYYY@{}}
\toprule
Method & Task focus & Morphology axis & Body configs \\
\midrule
DeepMimic~\cite{deepmimic} & Reference-guided skill imitation & None (single body) & 1 \\
PHC~\cite{phc} & Real-time humanoid imitation & None (single body) & 1 \\
PULSE~\cite{pulse} & Universal low-dimensional motion representation & None (single body) & 1 \\
OmniH2O~\cite{omnih2o} & Whole-body teleoperation & None (single physical robot) & 1 \\
ExBody2~\cite{exbody2} & Expressive whole-body control & None (single physical robot) & 1 \\
GMT~\cite{gmt} & Broad motion tracking (mixture-of-experts) & None (single body) & 1 \\
Won \& Lee~\cite{won2019} & Parametric body-shape control & Continuous (height, weight, proportions) & Parametric range; count not specified \\
Shared modular policies~\cite{smp2020} & Cross-topology transfer & Topology (limb count/structure), not anthropometric shape & N/A -- topology axis \\
\midrule
This work & Multi-morphology motion imitation & Continuous anthropometric shape (SMPL $\boldsymbol{\beta}$), matched simulation assets & \textbf{128}, over 20,951 paired clips \\
\bottomrule
\end{tabularx}
\end{table}

Kinematic retargeting and dynamic tracking address different requirements.
SMPL provides a common parameterized body representation~\cite{smpl}, and
SMPLSim constructs simulation-ready assets from these
models~\cite{smplsim}. A shared skeleton makes poses comparable, but does not
make global positions or required actuation independent of shape. Matching a
reference to its intended simulation asset is consequently part of the task
definition, not an optional evaluation convenience.

\subsection{Text-motion data, shape-conditioned generation, and refinement}

AMASS consolidates motion capture in a common body-model
representation~\cite{amass}. HumanML3D associates processed motion segments
with natural-language descriptions~\cite{humanml3d}. We inherit those source
identities and descriptions rather than claiming new motion capture or
shape-specific language annotation.

HUMOS generates identity-conditioned human motion using a cycle-consistent
learning formulation~\cite{humos}. We use its pretrained model to obtain
shape-specific references rather than as the physics controller or a
text-conditioned policy. The resulting sequences still require conversion to
the simulator's skeleton and coordinate conventions, and their dynamic
feasibility must be assessed in that simulator; our refinement is a
deliberately limited kinematic post-process, not a new generative model or a
full trajectory optimizer.

The controller uses standard Transformer operations~\cite{attention}, with
learned slot and type embeddings identifying temporal positions and
observation sources. An ablation also investigates actor-only morphology
modulation with adaLN-Zero, adapting the zero-initialized conditioning design
used in Diffusion Transformers~\cite{dit} as a conditioning mechanism rather
than for diffusion-based action generation.
